\section{Testing, Verification and Failover}\label{sec:testing}
% Rubric criterion 6 (15 %): CLI outputs (show ip route, traceroute, ping), failover scenarios documented and proven.
% Only measured facts here: a failover claimed in the report must be repeatable in the live demo.

\subsection{Verification method}\label{sec:test-method}
We verified the design at four levels. Each level proves something the previous one cannot, and none of the
first three replaces the lab.

\begin{table}[H]\centering\footnotesize
\caption{Verification levels}\label{tab:verification}
\begin{tabularx}{\textwidth}{@{}>{\raggedright\arraybackslash}p{2.6cm} >{\raggedright\arraybackslash}p{3.1cm} Y Y@{}}
\toprule
\textbf{Level} & \textbf{Tool} & \textbf{Proves} & \textbf{Cannot prove} \\
\midrule
1. Data model & \FactTestsModel{} unit tests & The plan is consistent: no overlapping subnets, capacity covers demand, unique addresses, root bridge = HSRP owner & That the rendered text means the same \\
2. Rendered configs & \FactTestsConfigs{} unit tests reading only the \texttt{.cfg} files & Both ends of every BGP session, OSPF adjacency, HSRP group, tunnel and trunk agree; every referenced policy exists & That the protocols converge to the intended paths \\
3. Routing emulation & Real OSPF/BGP daemons (FRRouting) in Linux namespaces & The routing logic converges and fails over (\cref{sec:emulation}) & IOS syntax, IPsec, Layer~2, IPv6, ACL/QoS behaviour \\
4. Lab & \todo{EVE-NG/GNS3}, IOSv/IOSvL2 & Everything, on the graded platform & -- \\
\bottomrule
\end{tabularx}
\end{table}

In the lab, the verification commands of every step are run from the automation host by
\path{Lab/collect_evidence.py}, which stores each output with the device prompt, the command and a
timestamp; the listings below are those files, unedited. Tests run on the VPCS hosts (ping, traceroute,
DHCP) are shown as screenshots.
\todo{note which commands, if any, the IOSv images rejected and were removed}

\subsection{Layer 2}\label{sec:test-l2}
\Cref{ev:s3_3} shows Po1 in use (\texttt{SU}) with both members bundled (\texttt{P}) and LACP as protocol.
\Cref{ev:s3_4} confirms that \dev{HCMC-CORE-SW01} is the root bridge for its VLAN group, and \cref{ev:s3_6}
shows the expected result on an access switch: one uplink forwarding as root port, the other in the
alternate (blocking) role.
\evidence[ev:s3_3]{s3_3}{LACP port-channel between the core switches}
\evidence[ev:s3_4]{s3_4}{Spanning-tree root placement on \dev{HCMC-CORE-SW01}}
\evidence[ev:s3_6]{s3_6}{Access switch: forwarding and alternate uplink}

\subsection{Gateway redundancy and IP services}\label{sec:test-fhrp}
The two HQ cores show complementary HSRP states -- each active for the VLANs of which it is the root bridge
and standby for the others (\cref{ev:s4_1a,ev:s4_1b}). The central DHCP server holds leases from all three
sites (\cref{ev:s4_2}), which proves that the relays and the return path through the tunnels work.
\evidence[ev:s4_1a]{s4_1a}{HSRP state on \dev{HCMC-CORE-SW01}}
\evidence[ev:s4_1b]{s4_1b}{HSRP state on \dev{HCMC-CORE-SW02}}
\evidence[ev:s4_2]{s4_2}{Central DHCP server: leases from all three sites}
\screenshot[fig:vpcs-dhcp]{vpcs_dhcp_hanoi.png}{A Hanoi host obtains its address from the HQ server; its gateway is the HSRP virtual address}{0.7\textwidth}

\subsection{Underlay}\label{sec:test-underlay}
At HQ the core switch has full adjacencies with the other core and both edge routers (\cref{ev:s5_1}). Its
routing table (\cref{ev:s5_2}) contains the default route as \texttt{O*E1} via \dev{HCMC-EDGE-RTR01}, the
router with the lower seed metric, and all loopbacks of the site are reachable (\cref{ev:s5_4}).
\evidence[ev:s5_1]{s5_1}{OSPF adjacencies at HQ}
\evidence[ev:s5_2]{s5_2}{OSPF routes including the default from the edge}
\evidence[ev:s5_4]{s5_4}{All HQ loopbacks reachable}

\subsection{Overlay and WAN}\label{sec:test-overlay}
On the HQ hub every BGP session is established -- the provider, the other HQ edge router and both branch
routers (\cref{ev:s6_1a}) -- and IKE shows \texttt{QM\_IDLE} security associations with growing packet
counters (\cref{ev:s6_2a}). The BGP table (\cref{ev:s6_3a}) contains the default route from the provider
and the aggregates of Hanoi and Da Nang with local preference 200. Hanoi learns the HQ and Da Nang aggregates
over its tunnel (\cref{ev:s6_4}), and end-to-end tests across the overlay succeed (\cref{ev:s6_7}).
\evidence[ev:s6_1a]{s6_1a}{BGP sessions on \dev{HCMC-EDGE-RTR01}}
\evidence[ev:s6_2a]{s6_2a}{IKE and IPsec security associations on \dev{HCMC-EDGE-RTR01}}
\evidence[ev:s6_3a]{s6_3a}{BGP table on \dev{HCMC-EDGE-RTR01}}
\evidence[ev:s6_4]{s6_4}{Hanoi learns the HQ and Da Nang aggregates over the tunnel}
\evidence[ev:s6_7]{s6_7}{End-to-end reachability and path across the overlay}

\subsection{Security}\label{sec:test-security}
\Cref{tab:acl-tests} lists the segmentation tests, run from the VPCS hosts. The blocked cases are as
important as the permitted ones: they prove the policy. The ACL hit counters (\cref{ev:s7_1}) show the deny
entries matching.

\begin{table}[H]\centering\footnotesize
\caption{Segmentation tests}\label{tab:acl-tests}
\begin{tabularx}{\textwidth}{@{}l >{\raggedright\arraybackslash}p{5.2cm} l l Y@{}}
\toprule
\textbf{\#} & \textbf{Test} & \textbf{Expected} & \textbf{Observed} & \textbf{Reason} \\
\midrule
1 & Executive PC $\rightarrow$ R\&D PC (ping) & reply & \todo{} & executives may initiate \\
2 & R\&D PC $\rightarrow$ executive PC (ping) & blocked & \todo{} & executive VLAN accepts replies only \\
3 & Guest PC $\rightarrow$ 8.8.8.8 & reply & \todo{} & Internet only \\
4 & Guest PC $\rightarrow$ DC server 10.1.100.10 & blocked & \todo{} & private destinations denied \\
5 & Hanoi R\&D PC $\rightarrow$ 8.8.8.8 & reply, NAT at Hanoi & \todo{} & local breakout \\
6 & Da Nang IoT $\rightarrow$ DC server (ping) & reply & \todo{} & infrastructure diagnostics allowed \\
7 & Da Nang IoT $\rightarrow$ R\&D PC / 8.8.8.8 & blocked & \todo{} & no lateral movement, no Internet \\
8 & Hanoi R\&D PC $\rightarrow$ automation host (VLAN 99) & blocked & \todo{} & management protected \\
9 & Automation host $\rightarrow$ SSH to a Hanoi router & login & \todo{} & VTY ACL allows management VLANs \\
\bottomrule
\end{tabularx}
\end{table}

\evidence[ev:s7_1]{s7_1}{Segmentation ACL hit counters at HQ}

Local breakout is visible in the NAT tables: traffic from an HQ host is translated at the HQ edge
(\cref{ev:s7_13}), traffic from a Hanoi host at the Hanoi edge, and the HQ edge shows no Hanoi translation
(\cref{ev:s7_13b}). On the access switches, DHCP snooping has a binding for every host (\cref{ev:s7_15}) and DAI
is active on the user VLANs (\cref{ev:s7_16}).
\evidence[ev:s7_13]{s7_13}{NAT overload at the HQ edge}
\evidence[ev:s7_13b]{s7_13b}{NAT at the branch: local Internet breakout}
\evidence[ev:s7_15]{s7_15}{DHCP snooping bindings}
\evidence[ev:s7_16]{s7_16}{Dynamic ARP inspection}

\subsection{Quality of service}\label{sec:test-qos}
Before the test (\cref{ev:s8_1}) the class counters of the WAN policy are near zero. After 300 pings marked
EF (\texttt{tos 184}) across the tunnel, the voice class has counted them in the priority queue
(\cref{ev:s8_2}). \todo{optional: result of the congestion test -- voice without drops while the default
class drops}
\evidence[ev:s8_1]{s8_1}{WAN policy counters before the test traffic}
\evidence[ev:s8_2]{s8_2}{Voice-marked traffic in the priority queue}

\subsection{Automation}\label{sec:test-automation}
\todo{a dry run; a real run adding VLAN 110 to the core and access switches with its summary; the listing of
\texttt{backups/}; the tail of \texttt{automation.log} including one deliberately caused error (wrong password or
unreachable device); \texttt{show vlan brief} with the new VLAN}
\screenshot[fig:auto-run]{automation_run.png}{Automation run: backup, push, verification and summary}{\textwidth}
\screenshot[fig:auto-log]{automation_log.png}{Automation log with a successful run and a handled connection error}{\textwidth}

\subsection{End-to-end connectivity and failover}\label{sec:failover}
For each scenario a host pings a host at another site continuously while the failure is injected; the
outage is the number of lost replies multiplied by the ping interval. The \emph{bound} column is what the
configured timers guarantee (\cref{sec:fhrp,sec:underlay,sec:overlay}).

\begin{table}[H]\centering\small
\caption{Failover scenarios in the lab}\label{tab:failover}
\begin{tabularx}{\textwidth}{@{}l Y Y l l@{}}
\toprule
\textbf{\#} & \textbf{Failure injected} & \textbf{Mechanism} & \textbf{Bound} & \textbf{Measured} \\
\midrule
1 & One member of Po1 down & LACP removes the member, no topology change & $<1$\,s & \todo{} \\
2 & \dev{HCMC-CORE-SW01} down & HSRP hold time, RSTP, OSPF dead interval & $\leq 4$\,s & \todo{} \\
3 & ISP-A uplink of HQ down & default withdrawn from OSPF; tunnels 2/4 & $\leq 9$\,s & \todo{} \\
4 & ISP-A uplink of Hanoi down & BGP hold time; route tracking moves HSRP & $\leq 10$\,s & \todo{} \\
5 & \dev{HCMC-EDGE-RTR01} down & OSPF dead interval, iBGP hold time & $\leq 9$\,s & \todo{} \\
\bottomrule
\end{tabularx}
\end{table}

\evidence[ev:s10_2]{s10_2_core1_down}{\dev{HCMC-CORE-SW02} after the failure of core 1: active for every VLAN}

\subsection{Supplementary: routing emulation with FRRouting}\label{sec:emulation}
Before building the lab we executed the routing logic of the configurations. A translator converts the IPv4
routing sections of the generated IOS configurations (interfaces, OSPF, BGP, prefix lists, route maps,
static routes, NAT) line by line into FRRouting \cite{frr} configurations, and a test harness runs one OSPF/BGP
instance per router in Linux network namespaces. Every link passes through a software bridge, as in EVE-NG,
so a failed neighbour leaves the far end up and must be detected by protocol timers. HSRP is replaced by a
rule engine that applies the configured priorities, tracking, preemption and hold times. The emulation does
not cover IOS syntax, IPsec (tunnels are unencrypted), Layer~2, IPv6, ACLs or QoS, and its outputs are FRR
outputs, not Cisco outputs.

In the baseline all \FactEmuBgpSessions{} BGP sessions were established and all eight test flows -- between
the sites and towards the Internet -- were reachable. \Cref{tab:emulation} lists the injected failures:
\FactEmuPassed{} of the eight single failures recovered with no flow left broken, the worst outages (about
9\,s) occurring where only the BGP hold timer can detect the failure. Scenario S9 is a deliberate double
failure that removes the edge-to-edge link; it reproduces the routing loop that this link prevents
(\cref{sec:challenges}).

\IfFileExists{tables/emulation.tex}{% GENERATED by Tools/gen_latex.py - do not edit; change Tools/tvi_model.py and regenerate.
\begin{table}[htbp]
\centering\footnotesize
\caption{FRR emulation of the routing design: failure scenarios (supplementary evidence, not the IOSv lab)}\label{tab:emulation}
\begin{tabularx}{\textwidth}{l X l r r}
\toprule
\textbf{ID} & \textbf{Injected failure} & \textbf{Result} & \textbf{Worst outage} & \textbf{Flows down at end} \\
\midrule
S1 & HQ ISP-A uplink fails (HCMC-EDGE-RTR01 Gi0/0) & pass & 8.8\,s & 0/8 \\
S2 & Hanoi ISP-A uplink fails (HANOI-EDGE-RTR01 Gi0/0) & pass & 8.0\,s & 0/8 \\
S3 & HQ edge router 1 fails completely (HCMC-EDGE-RTR01) & pass & 8.8\,s & 0/8 \\
S4 & HQ core switch 1 fails completely (HCMC-CORE-SW01) & pass & 3.8\,s & 0/8 \\
S5 & Core1-Edge1 routed link fails (HCMC-CORE-SW01 Gi1/0) & pass & 1.0\,s & 0/8 \\
S6 & Da Nang ISP-B uplink fails (DANANG-EDGE-RTR02 Gi0/0) & pass & 1.2\,s & 0/8 \\
S7 & Hanoi edge router 2 fails completely (HANOI-EDGE-RTR02) & pass & 1.8\,s & 0/8 \\
S8 & Whole provider ISP-A fails (ISP-A-RTR) & pass & 8.6\,s & 0/8 \\
S9 & DOUBLE failure: HQ edge-to-edge link + Hanoi ISP-A uplink & \textbf{fail} & 26.0\,s & 4/8 \\
\bottomrule
\end{tabularx}
\end{table}
}{}
